% tex/30_rankcat.tex
\section{Categories: the canonical \texorpdfstring{$\le$}{Le}-lane}
\label{sec:rankcat}

This section packages rank morphisms into a category and fixes the lane story.

\subsection{Morphisms in the \texorpdfstring{$\le$}{Le}-lane}
\begin{definition}[\leaninline{RankHomLe}]
A \emph{\(\le\)-morphism} \(f : (X,\rank_X)\to (Y,\rank_Y)\) consists of:
\begin{itemize}
\item a map \(f.\mathrm{map}:X\to Y\),
\item an index map \(f.\mathrm{indexMap}:\alpha\to\beta\),
\item monotonicity of \(f.\mathrm{indexMap}\),
\item the inequality \(\rank_Y(f(x)) \le f.\mathrm{indexMap}(\rank_X(x))\).
\end{itemize}
Composition is designed to use only transitivity and monotonicity.
\end{definition}

\subsection{The category \leaninline{RankCatLe}}
\begin{definition}[\leaninline{RankCatLe α}]
Objects are pairs \(\Sigma X, \Ranked\,\alpha\,X\) and morphisms are \leaninline{RankHomLe}.
\end{definition}

\subsection{Lane bridge: Eq \texorpdfstring{$\to$}{->} Le \texorpdfstring{$\to$}{->} D}
We support weakening morphisms:
\[
  \HomEq \longrightarrow \HomLe \longrightarrow \HomD.
\]
Keep this diagrammatic and implementation-oriented (not heavy category theory).

\subsection{Where to paste T3}
Paste (or paraphrase) the frozen T3 text here:
\begin{itemize}
\item \leaninline{RankHomLe}: definition + \leaninline{id}/\leaninline{comp} and the “stable proof” story;
\item \leaninline{RankCatLe}: category instance;
\item bridge file \leaninline{RankBridge_EqLeD.lean}: the Eq→Le→D adapters.
\end{itemize}
